\documentclass[11pt]{article}

\usepackage{amsmath, amssymb}
\usepackage{booktabs}
\usepackage{geometry}
\geometry{margin=1in}

\title{Lecture 18 Scribe \\ \large Marginal PMF, Correlation, Covariance, and Joint Gaussian Random Variables}
\date{}
\author{}

\begin{document}

\maketitle

\section*{1. Marginal Probability Mass Function (Marginal PMF)}

\subsection*{Definition}

For discrete random variables $X$ and $Y$, the joint PMF is

\[
p_{XY}(x,y) = P(X=x, Y=y)
\]

The marginal PMF of a variable is obtained by summing the joint PMF over all possible values of the other variable.

\[
p_X(x) = \sum_y p_{XY}(x,y)
\]

\[
p_Y(y) = \sum_x p_{XY}(x,y)
\]

\textbf{Key Idea:}

The marginal PMF of one variable is obtained by summing the joint PMF over all possible values of the other variable.

\subsection*{Marginal PMF from a Joint PMF Table}

Given a joint PMF table $p_{XY}(x,y)$:

\begin{itemize}
\item Row sums give $p_X(x)$
\item Column sums give $p_Y(y)$
\end{itemize}

\[
p_X(x_i) = \sum_j p_{ij}
\]

\[
p_Y(y_j) = \sum_i p_{ij}
\]

Normalization conditions:

\[
\sum_i p_X(x_i) = 1
\]

\[
\sum_j p_Y(y_j) = 1
\]

\subsection*{Example}

Joint PMF table:

\begin{center}
\begin{tabular}{c|cc}
$p_{XY}(x,y)$ & $y=0$ & $y=1$ \\
\midrule
$x=0$ & $1/8$ & $1/4$ \\
$x=1$ & $3/8$ & $1/4$
\end{tabular}
\end{center}

Compute marginal PMFs

\[
p_X(0) = \frac{1}{8} + \frac{1}{4} = \frac{3}{8}
\]

\[
p_X(1) = \frac{3}{8} + \frac{1}{4} = \frac{5}{8}
\]

\[
p_Y(0) = \frac{1}{8} + \frac{3}{8} = \frac{1}{2}
\]

\[
p_Y(1) = \frac{1}{4} + \frac{1}{4} = \frac{1}{2}
\]

Thus,

\[
p_X(x)=
\begin{cases}
3/8, & x=0 \\
5/8, & x=1
\end{cases}
\]

\[
p_Y(y)=
\begin{cases}
1/2, & y=0 \\
1/2, & y=1
\end{cases}
\]

\section*{2. Correlation}

\subsection*{Definition}

A basic way to measure joint behavior between two random variables is the product moment.

\[
R_{XY} = E[XY]
\]

This represents the average value of the product $XY$ and is often called a raw correlation measure.

\subsection*{Interpretation}

\begin{itemize}
\item $R_{XY}$ depends on the scale of $X$ and $Y$
\item Computed about the origin, not about the means
\item Mixes mean values, spread, and dependence
\end{itemize}

If

\[
R_{XY}=0
\]

then $X$ and $Y$ are orthogonal about the origin.

To remove the effect of means, the variables are centered, which leads to covariance.

\section*{3. Covariance}

\subsection*{Definition}

Covariance measures how two variables vary together after removing their means.

\[
\text{Cov}(X,Y) = E[(X-\mu_X)(Y-\mu_Y)]
\]

where

\[
\mu_X = E[X], \quad \mu_Y = E[Y]
\]

\subsection*{Interpretation}

\textbf{Positive covariance}

\[
\text{Cov}(X,Y) > 0
\]

Deviations from the means tend to have the same sign.

\textbf{Negative covariance}

\[
\text{Cov}(X,Y) < 0
\]

When one variable increases, the other tends to decrease.

\textbf{Zero covariance}

\[
\text{Cov}(X,Y) = 0
\]

There is no linear co-variation between variables.

\subsection*{Algebraic Form}

Starting from

\[
\text{Cov}(X,Y) = E[(X-\mu_X)(Y-\mu_Y)]
\]

Expanding

\[
= E[XY - X\mu_Y - Y\mu_X + \mu_X\mu_Y]
\]

\[
= E[XY] - \mu_YE[X] - \mu_XE[Y] + \mu_X\mu_Y
\]

Since

\[
\mu_X = E[X], \quad \mu_Y = E[Y]
\]

\[
\text{Cov}(X,Y) = E[XY] - \mu_X\mu_Y
\]

or

\[
\text{Cov}(X,Y) = E[XY] - E[X]E[Y]
\]

\subsection*{Properties}

\begin{itemize}
\item Symmetry
\[
\text{Cov}(X,Y) = \text{Cov}(Y,X)
\]

\item Variance as special case
\[
\text{Cov}(X,X) = \text{Var}(X)
\]

\item Linearity with constants
\[
\text{Cov}(aX+b, cY+d) = ac\,\text{Cov}(X,Y)
\]

\item Independence implies zero covariance
\end{itemize}

If $X \perp Y$

\[
\text{Cov}(X,Y)=0
\]

However, zero covariance does not necessarily imply independence.

\subsection*{Example}

If

\[
Y = -2X + 3
\]

Find $\text{Cov}(X,Y)$.

\[
XY = X(-2X+3) = -2X^2 + 3X
\]

\[
E[Y] = E[-2X+3] = -2E[X] + 3
\]

\[
\text{Cov}(X,Y) = -2E[X^2] + 2(E[X])^2
\]

Using

\[
\text{Var}(X) = E[X^2] - (E[X])^2
\]

\[
\text{Cov}(X,Y) = -2\text{Var}(X)
\]

\section*{4. Pearson Correlation Coefficient}

\[
\rho_{XY} =
\frac{\text{Cov}(X,Y)}{\sqrt{\text{Var}(X)\text{Var}(Y)}}
\]

or

\[
\rho_{XY} =
\frac{\text{Cov}(X,Y)}{\sigma_X\sigma_Y}
\]

Properties

\begin{itemize}
\item Normalized version of covariance
\item Unit-free
\item Measures strength and direction of linear relationship
\end{itemize}

Range

\[
-1 \le \rho_{XY} \le 1
\]

Interpretation

\begin{itemize}
\item $\rho_{XY} > 0$ : positive linear trend
\item $\rho_{XY} < 0$ : negative linear trend
\item $\rho_{XY} = 0$ : no linear correlation
\end{itemize}

\section*{5. Example (Joint PDF)}

Given

\[
f_{XY}(x,y)=
\begin{cases}
\dfrac{xy}{96}, & 0<x<4,\; 1<y<5 \\
0, & \text{otherwise}
\end{cases}
\]

Results

\[
E[X] = \frac{8}{3}
\]

\[
E[Y] = \frac{31}{9}
\]

\[
E[XY] = \frac{248}{27}
\]

\[
E[2X+3Y] = \frac{47}{3}
\]

\[
Var(X) = \frac{8}{9}
\]

\[
Var(Y) = \frac{92}{81}
\]

\[
Cov(X,Y)=0
\]

\[
\rho_{XY}=0
\]

\section*{6. Jointly Gaussian Random Variables}

\subsection*{Definition}

A pair $(X,Y)$ is jointly Gaussian if

\begin{itemize}
\item The joint density follows a bivariate Gaussian distribution
\item The marginal distributions are Gaussian
\end{itemize}

\subsection*{Joint Gaussian PDF}

\[
f_{XY}(x,y)=
\frac{1}{2\pi\sigma_X\sigma_Y\sqrt{1-\rho^2}}
\exp
\left(
-\frac{1}{2(1-\rho^2)}
\left[
\left(\frac{x-\mu_X}{\sigma_X}\right)^2
+
\left(\frac{y-\mu_Y}{\sigma_Y}\right)^2
-
2\rho
\left(\frac{x-\mu_X}{\sigma_X}\right)
\left(\frac{y-\mu_Y}{\sigma_Y}\right)
\right]
\right)
\]

where

\begin{itemize}
\item $\mu_X,\mu_Y$ : means
\item $\sigma_X,\sigma_Y$ : standard deviations
\item $\rho$ : Pearson correlation coefficient
\end{itemize}

\subsection*{Real-Life Interpretation}

Example variables

\begin{itemize}
\item $X$ : Atmospheric CO$_2$ concentration
\item $Y$ : Global temperature anomaly
\end{itemize}

If $\rho > 0$, higher CO$_2$ levels tend to occur with higher temperatures.

\end{document}